\section{Introduction}

Small aerial robots rarely fly the vehicle they were tuned on. A delivery adds payload, a propeller chips, the wind picks up and the companion computer adds latency. Classical control handles part of this: a geometric tracking controller~\cite{lee2010geometric} with integral action rejects slow disturbances, and L1 adaptive augmentation estimates and cancels a lumped disturbance online~\cite{wu2022l1quad}. Learning-based adaptation promises more: meta-learning trains a model that adapts from a few samples of a new task~\cite{finn2017maml}, and learned adaptation has produced agile flight in strong winds~\cite{oconnell2022neuralfly} and single controllers for very different quadrotors~\cite{zhang2022nearhover,eschmann2025raptor}.

Two questions remain open. First, meta-reinforcement-learning algorithms are typically evaluated against each other on simulated robot benchmarks such as locomotion and manipulation suites~\cite{yu2019metaworld,beck2025tutorial}, rather than against the classical controllers and planners that would run in the same loop. Second, a deployed autonomy stack has several places where learning could enter (the low-level controller, its gains, the local planner), and it is not clear which of them benefits from meta-learning.

We study both questions in one simulation suite. A batched quadrotor model parameterised after a Holybro X500 V2 runs a classical stack: a geometric controller with collective-thrust and body-rate (CTBR) output, a simulated RealSense D435i depth camera, a log-odds occupancy map and an A* planner (Fig.~\ref{fig:pipeline}). Meta-learning plugs into three places: a residual on the controller, the controller's gains, and the local planner. Each must adapt to a new task drawn from payload, motor wear, drag, wind, gusts and latency.

Our contributions are:
\begin{itemize}
\item an open simulation suite with three classical+meta combinations and a reproducible benchmark protocol (Secs.~\ref{sec:problem}--\ref{sec:protocol});
\item a comparison of seven meta-learning algorithms with nominal, L1-adaptive, domain-randomised (DR) and fine-tuning baselines on held-out and out-of-distribution (OOD) tasks (Sec.~\ref{sec:results});
\item four findings: a learned residual beats tuned classical adaptation; meta-learning helps most when it adapts classical gains; gradient-based adaptation at the meta-trained step size degrades learned navigation planners, and the classical planner stays the most reliable on the full stack; and no method survives the OOD tasks;
\item a hardware recommendation and staged sim-to-real plan for testing the stack (Sec.~\ref{sec:platform}).
\end{itemize}
Code, results and a replay of recorded flights are available at \url{https://github.com/amalmehta/DroneAutonomy}.

\begin{figure*}[t]
\centering
% Fig. 1: the stack and the three places meta-learning plugs in (hand-drawn TikZ, no data)
\begin{tikzpicture}[
  box/.style={draw, rounded corners=1.5pt, align=center, font=\footnotesize, minimum height=8.5mm, inner sep=2pt, text width=16mm},
  learn/.style={box, draw=orange!80!black, very thick, fill=orange!8},
  arr/.style={-{Stealth[length=2mm]}, thick},
  tag/.style={font=\scriptsize\itshape, text=orange!70!black, align=center}]
  \node[box] (cam) {Depth camera\\(D435i model)};
  \node[box, right=4mm of cam] (map) {Occupancy map\\(log-odds voxels)};
  \node[box, right=4mm of map] (astar) {A* planner\\+ carrot};
  \node[learn, right=4mm of astar] (local) {Local planner\\(velocity)};
  \node[learn, right=4mm of local] (ctrl) {Geometric\\controller};
  \node[box, right=9mm of ctrl] (dyn) {Quadrotor\\dynamics};
  \draw[arr] (cam) -- (map); \draw[arr] (map) -- (astar); \draw[arr] (astar) -- (local);
  \draw[arr] (local) -- (ctrl);
  \draw[arr] (ctrl) -- node[above, font=\scriptsize] {CTBR} (dyn);
  \draw[arr] (dyn.south) -- ++(0,-4mm) -| node[pos=0.25, below, font=\scriptsize] {pose, depth rays} (cam.south);
  \node[tag, above=1mm of local] {Combination 3:\\meta-RL planner};
  \node[tag, above=1mm of ctrl] {Combinations 1, 2:\\residual or gains};
  \node[font=\scriptsize, align=center, below=7mm of ctrl] (task) {Task: payload, motor wear,\\drag, wind, gusts, latency};
  \draw[arr, dashed] (task.east) -| (dyn.south east);
\end{tikzpicture}

\caption{Meta-learning plugs into three places of one classical stack. The local planner (Combination 3) turns depth and a carrot point on the A* path into a velocity command. The geometric controller turns that into collective thrust and body rates, with a learned residual (Combination 1) or meta-learned gains (Combination 2). The task perturbs the dynamics.}
\label{fig:pipeline}
\end{figure*}
